% ---------------------------------------------------------------------------
% Shared preamble of every ClaimFlow sprint report. Build with XeLaTeX from
% the sprint folder (bun run pdf sprint-N). Before \input{../preamble},
% main.tex defines \sprintname (Sprint 1) and \deliverable (Project
% Deliverable 1).
%
% Fonts: Times New Roman and Courier New ship with Windows and macOS and both
% carry the arrows and comparison signs the sources use. A glyph neither has
% shows as "Missing character" in main.log; check-report counts those.
% ---------------------------------------------------------------------------
\usepackage{fontspec}
\setmainfont{Times New Roman}
\setmonofont{Courier New}

% Student names and IDs, from docs/reports/.env through bun run identity.
\InputIfFileExists{build/identity.tex}{}{%
  \errmessage{build/identity.tex is missing: run bun run identity <sprint>}}

\usepackage[a4paper, margin=2cm]{geometry}
\usepackage{graphicx}
\usepackage{pdflscape}
\usepackage{booktabs}
\usepackage{longtable}
\usepackage{array}
\usepackage{xcolor}
\usepackage{fvextra}
\usepackage{float}
\usepackage{caption}
\usepackage{etoolbox}
\usepackage{refcount}
\usepackage{seqsplit}
\usepackage{fancyhdr}

% Colours before hyperref, which reads linkcolor as it loads.
\definecolor{brand}{HTML}{1B5E3F}
\definecolor{todored}{HTML}{C62828}
\usepackage[colorlinks=true, urlcolor=blue!55!black, linkcolor=brand,
            pdftitle={ClaimFlow \sprintname\space report},
            pdfauthor={\studentOneName; \studentTwoName}]{hyperref}

\setlength{\parindent}{0pt}
\setlength{\parskip}{5pt}
\setlength{\emergencystretch}{3em}
\hfuzz=3pt % ignore overfull boxes under 3 pt, which no reader can see
\captionsetup{font=small}
\setcounter{tocdepth}{2}
\setcounter{secnumdepth}{0}

\pagestyle{fancy}
\fancyhf{}
\fancyfoot[L]{\footnotesize ClaimFlow, \sprintname\space report}
\fancyfoot[R]{\footnotesize page \thepage}
\renewcommand{\headrulewidth}{0pt}
\fancypagestyle{plain}{\fancyhf{}\fancyfoot[R]{\footnotesize page \thepage}%
  \renewcommand{\headrulewidth}{0pt}}

% ---- Evidence that is not there yet ---------------------------------------

% Red inline text. check-report counts every \todo and \pending.
\newcommand{\todo}[1]{\textcolor{todored}{\textbf{[TODO: #1]}}}

% A red box where a capture or section is still waiting for something.
% \pending{what it shows}{what it waits for}
\newcommand{\pending}[2]{%
  \par\medskip\noindent\fcolorbox{todored}{white}{\parbox{0.97\linewidth}{%
    \textcolor{todored}{\textbf{[PENDING: #1]}}\\
    \small\itshape Waits for: #2}}\par\medskip}

% ---- Parts -----------------------------------------------------------------

% \reportpart{A}{Title}: starts a page, adds the contents entry, and sets
% part:A:start. \reportpartend{A} sets part:A:end on the last page, so the
% navigation table can print real page ranges with \pageref.
\newcommand{\reportpart}[2]{%
  \clearpage
  \phantomsection
  \addcontentsline{toc}{section}{Part #1: #2}%
  \label{part:#1:start}%
  {\noindent\bfseries\Large\color{brand} Part #1: #2\par}%
  \vspace{2pt}\noindent\textcolor{brand}{\rule{\linewidth}{1.2pt}}\par
  \vspace{6pt}}
\newcommand{\reportpartend}[1]{\par\label{part:#1:end}}
% \pages{A} prints "pp. 5--12" or "p. 5"
\newcommand{\pages}[1]{%
  \ifnumequal{\getpagerefnumber{part:#1:start}}{\getpagerefnumber{part:#1:end}}%
    {p.\ \pageref{part:#1:start}}%
    {pp.\ \pageref{part:#1:start}--\pageref{part:#1:end}}}

% Sub-headings inside a part.
\newcommand{\partsection}[2]{%
  \phantomsection\addcontentsline{toc}{subsection}{#1\ \ #2}%
  \par\bigskip{\noindent\bfseries\large\color{brand} #1\ \ #2\par}\smallskip}

% ---- Evidence figures --------------------------------------------------------

\newcommand{\evidencemaxheight}{22cm}

% The red box drawn when a figure file does not exist, so the build never
% fails on a missing capture.
\newcommand{\figmissing}[2]{%
  \fbox{\parbox[c][3cm][c]{0.94\linewidth}{\centering
    \textcolor{todored}{\textbf{[PENDING FIGURE: \detokenize{#1}]}}\\[4pt]
    \small\itshape #2}}}

% \evidence{figures/x.png}{Caption}: one full-width numbered figure.
\newcommand{\evidence}[2]{%
  \IfFileExists{#1}%
    {\begin{figure}[H]\centering
       \fbox{\includegraphics[width=0.96\linewidth,height=\evidencemaxheight,keepaspectratio]{#1}}%
       \caption{#2}\end{figure}}%
    {\begin{figure}[H]\centering\figmissing{#1}{#2}\caption{#2}\end{figure}}}

% \evidencepart{figures/x-p2.png}{Caption}: one unnumbered piece of a sliced capture.
\newcommand{\evidencepart}[2]{%
  \IfFileExists{#1}%
    {\begin{figure}[H]\centering
       \fbox{\includegraphics[width=0.96\linewidth,height=\evidencemaxheight,keepaspectratio]{#1}}%
       \caption*{\small\itshape #2}\end{figure}}%
    {\begin{figure}[H]\centering\figmissing{#1}{#2}\caption*{\small\itshape #2}\end{figure}}}

% \evidencehalf{file}{caption}: wide capture, at most 11.5 cm tall, two to a page.
\newcommand{\evidencehalf}[2]{%
  \IfFileExists{#1}%
    {\begin{figure}[H]\centering
       \fbox{\includegraphics[width=0.96\linewidth,height=11.5cm,keepaspectratio]{#1}}%
       \caption{#2}\end{figure}}%
    {\begin{figure}[H]\centering\figmissing{#1}{#2}\caption{#2}\end{figure}}}

% \diagramfull{file}{caption}: one Mermaid diagram of a captured document,
% photographed on its own at 3x by scripts/capture-diagrams.ts, at full text width.
\newcommand{\diagramfull}[2]{%
  \IfFileExists{#1}%
    {\begin{figure}[H]\centering
       \fbox{\includegraphics[width=0.96\linewidth,height=\evidencemaxheight,keepaspectratio]{#1}}%
       \caption{#2}\end{figure}}%
    {\begin{figure}[H]\centering\figmissing{#1}{#2}\caption{#2}\end{figure}}}

% \diagramwide{file}{caption}: a diagram much wider than tall, on its own landscape page.
\newcommand{\diagramwide}[2]{%
  \IfFileExists{#1}%
    {\begin{landscape}\begin{figure}[H]\centering
       \fbox{\includegraphics[width=0.96\linewidth,height=15cm,keepaspectratio]{#1}}%
       \caption{#2}\end{figure}\end{landscape}}%
    {\begin{figure}[H]\centering\figmissing{#1}{#2}\caption{#2}\end{figure}}}

% \evidencesmall{file}{caption}{width fraction}: a capture that is small to begin with.
\newcommand{\evidencesmall}[3]{%
  \IfFileExists{#1}%
    {\begin{figure}[H]\centering
       \fbox{\includegraphics[width=#3\linewidth,keepaspectratio]{#1}}%
       \caption{#2}\end{figure}}%
    {\begin{figure}[H]\centering\figmissing{#1}{#2}\caption{#2}\end{figure}}}

% \docsource{docs/srs.md}{sprint-1}{abc1234}: the line above each captured document.
\newcommand{\docsource}[3]{%
  \par\noindent{\small\textbf{Source:} \texttt{\detokenize{#1}} as GitHub renders it,
  captured from \texttt{\detokenize{#2}} at commit \texttt{\detokenize{#3}}.}\par\smallskip}

% ---- Code and terminal output ------------------------------------------------

\DefineVerbatimEnvironment{shell}{Verbatim}
  {fontsize=\footnotesize, frame=single, framesep=5pt, rulecolor=\color{black!25},
   breaklines=true, breaksymbolleft={}}
\DefineVerbatimEnvironment{codeblock}{Verbatim}
  {fontsize=\scriptsize, frame=single, framesep=5pt, rulecolor=\color{black!25},
   breaklines=true, breaksymbolleft={}}

% ---- Title block --------------------------------------------------------------

\newcommand{\reporttitle}{%
  \begin{center}
    {\Large\bfseries CPE 334 --- Introduction to Software Engineering\par}
    \vspace{3pt}
    {\large \deliverable: \sprintname\space report\par}
    \vspace{10pt}
    {\Huge\bfseries\color{brand} ClaimFlow\par}
    {\large Theme 5, Department Expense Reimbursement System\par}
  \end{center}}
